\documentclass[11pt,a4paper]{article}

\input{../../sablony/pripravy/priprava.sty}

\nastavPripravu{4. ročník}{Kombinatorika}{39}{Variace s opakováním}

\begin{document}

\blok{Motivace a návaznost}

Čtyřmístný PIN má na každé pozici na výběr číslici od 0 do 9. Co se mění oproti permutacím s předem danými četnostmi?

\poznamka{Předpoklady: pravidlo součinu, variace bez opakování a permutace s opakováním. Cíl: rozlišit volné opakování symbolů od pevných četností a řešit omezené kódy.}

\blok{Výklad a odvození}

\begin{definiceBox}
\textbf{Variace s opakováním} délky $k$ z $n$ různých prvků je uspořádaná $k$-tice, v níž se smí každý prvek opakovat. Pro $n\ge1$, $k\ge0$ platí
\[
\overline{V}(n,k)=n^k.
\]
\end{definiceBox}
Kód smí začínat nulou, pokud zadání neurčí opak. Číslo začínat nulou nesmí. Při zákazu opakování sousedních znaků dostává první pozice $n$ možností, každá další $n-1$: celkem $n(n-1)^{k-1}$ pro $k\ge1$.

\blok{Společně řešené příklady}

\begin{prikladReseniBox}[title={Příklad 1}]
Kolik existuje čtyřmístných PINů z deseti číslic, pokud se číslice mohou opakovat?
\tcblower
$10^4=10000$. Počáteční nula je dovolena.
\end{prikladReseniBox}



\begin{prikladReseniBox}[title={Příklad 2}]
Kolik čtyřciferných čísel existuje, smějí-li se cifry opakovat?
\tcblower
První cifra $9$ možností, další po $10$: $9\cdot10^3=9000$.
\end{prikladReseniBox}



\begin{prikladReseniBox}[title={Příklad 3}]
Kolik pětiznakových kódů lze vytvořit z písmen A, B, C, pokud nesmějí být dva sousední znaky stejné?
\tcblower
První pozice tři možnosti, další vždy dvě: $3\cdot2^4=48$.
\end{prikladReseniBox}



\blok{Samostatná práce}

\begin{zadaniReseniBox}[title={Samostatná práce 1}]
Kolik trojmístných hesel z pěti písmen existuje s opakováním?
\tcblower
$5^3=125$.
\end{zadaniReseniBox}



\begin{zadaniReseniBox}[title={Samostatná práce 2}]
Kolik dvouciferných čísel se stejnými ciframi existuje?
\tcblower
Devět: $11,22,\ldots,99$.
\end{zadaniReseniBox}



\begin{zadaniReseniBox}[title={Samostatná práce 3}]
Kolik čtyřznakových kódů nad abecedou A, B, C obsahuje právě dvě A?
\tcblower
Polohy A vybereme $\binom42=6$ způsoby a na zbývající dvě místa po dvou možnostech B,C. Celkem $6\cdot2^2=24$.
\end{zadaniReseniBox}



\blok{Rozšíření}

\begin{zadaniReseniBox}[title={Rozšíření 1}]
Kolik čtyřmístných PINů obsahuje alespoň jednou číslici 7?
\tcblower
Všech PINů $10^4$, bez sedmičky $9^4$. Rozdíl $10000-6561=3439$.
\end{zadaniReseniBox}



\begin{zadaniReseniBox}[title={Rozšíření 2}]
\textbf{Volitelné:} Ze čtyř druhů zmrzliny vybereme tři kopečky, pořadí chutí nerozlišujeme a jednu chuť můžeme zvolit vícekrát. Kolik výběrů vznikne?
\tcblower
Jde o kombinace s opakováním. Počet je $\binom{4+3-1}{3}=\binom63=20$. Lze ověřit výpisem možností nebo modelem hvězdiček a přepážek. Nejde o povinný vzorec této hodiny.
\end{zadaniReseniBox}



\blok{Domácí úkol}

\poznamka{Doporučená práce doma: přibližně 15--30 minut. Úlohy jsou seřazeny od základních k náročnějším; učitel může poslední podúlohu označit jako volitelnou.}

\begin{zadaniReseniBox}[title={Domácí úkol}]

\begin{enumerate}[label=\alph*),leftmargin=*,itemsep=2pt]

\item Kolik je pětimístných PINů, pokud lze opakovat všechny číslice?
\item Kolik čtyřciferných čísel existuje při povoleném opakování číslic?
\item Kolik čtyřznakových hesel nad třemi písmeny A, B, C nemá dva sousední znaky stejné?
\item Kolik trojznakových kódů z číslic $0,1,2,3,4$ obsahuje alespoň jednu nulu?
\item Vysvětli rozdíl mezi permutací písmen slova AABB a čtyřznakovým kódem z písmen A, B s opakováním.
\item Kolik pětiznakových kódů nad písmeny A, B, C obsahuje právě jedno A?

\end{enumerate}

\tcblower

\begin{enumerate}[label=\alph*),leftmargin=*,itemsep=2pt]

\item $10^5=100000$.
\item $9\cdot10^3=9000$.
\item $3\cdot2^3=24$.
\item $5^3-4^3=125-64=61$.
\item Pevné dvě A a dvě B: $4!/(2!2!)=6$; volné četnosti: $2^4=16$.
\item Poloha A: pět možností, zbývající čtyři místa po dvou možnostech: $5\cdot2^4=80$.

\end{enumerate}

\end{zadaniReseniBox}

\poznamka{Orientační práce v hodině: 0--10 min kontrast kód a slovo; 10--24 min odvození; 24--38 min omezení; 38--45 min shrnutí.}

\poznamka{Kombinace s opakováním jsou pouze rozšiřující možnost; základní test by je neměl vyžadovat bez samostatného probrání.}

\end{document}
